\begin{frame}{자주 묻는 질문 --- PCA (3/3)}
  \small
  \begin{itemize}
    \item \kb{Q. 라벨을 모르는데, 분류 문제에도 쓸 수 있나?}
      PCA 자체는 라벨을 쓰지 않는다(비지도). 그러나 \textbf{전처리}로서
      자주 쓰임 --- ``PCA 축소 $\to$ 축소된 특징으로 분류기 학습''은
      표준. \; ``라벨을 구분하는 방향''은 못 찾지만, 중복/노이즈 축을
      정리하는 역할. \; 라벨을 아는 방법(지도 차원축소, 예: LDA)이
      필요하면 다른 기법으로.
    \item \kb{Q. 주성분의 방향을 ``해석''할 수 있나?}
      주성분은 원본 특징의 \textbf{선형 결합} --- 예:
      ``PC1 = 0.8$\times$꽃잎길이 + 0.6$\times$꽃잎너비 $-$ $\cdots$''.
      \; 가중치의 \textbf{크기와 부호}(이 가중치를 \kb{로딩}(loading)이라
      부름)으로 ``이 축이 어떤 원본 특징의 변동을 주로 담는지''를 판독.
      \par ``이 방향이 `꽃의 크기'다''처럼 의미를 부여하는 것은
      \textbf{해석자의 몫} --- PCA가 자동으로 ``의미''를 알려주지는
      않는다.
  \end{itemize}
\end{frame}
